\documentclass[10pt,a4paper]{amsart}
\usepackage[margin=19mm]{geometry}
\usepackage{amsmath,amssymb,mathtools}
\usepackage[scr=rsfs,cal=euler]{mathalfa}
\usepackage{xfrac}
\usepackage[unicode,hidelinks,bookmarks=false]{hyperref}
\newcommand{\ie}{i.e.\ }
\newcommand{\cf}{cf.\ }
\newcommand{\resp}{respectively\ }
\newcommand{\Subsection}{\subsection}
\providecommand{\nobreakdash}{\nobreak\hbox{-}}
\setlength{\emergencystretch}{2em}
\multlinegap0pt
\usepackage{bm}
\usepackage{bbm}
\usepackage{dsfont}
\usepackage{xspace}
\usepackage{cancel}
\usepackage{mathtools}
\usepackage{amsthm}
\makeatletter
\@ifundefined{theorem}{\newtheorem{theorem}{Theorem}}{}
\@ifundefined{lemma}{\newtheorem{lemma}[theorem]{Lemma}}{}
\makeatother
\newcommand{\F}{{\mathbb{F}}}
\title[Constructions of LCPs and LCD codes from twisted Reed-Solomo]{Constructions of LCPs and LCD codes from twisted Reed-Solomon codes}
\author{Shuo Sun, Wenwen Chen, Chao Liu, Yaozong Zhang, Xiaoqiang Wang}
\date{Source: arXiv:2609.16921v1 (CC BY 4.0)}
\begin{document}
\maketitle

\noindent\emph{Source and licence.} Constructions of LCPs and LCD codes from twisted Reed-Solomon codes. Version arXiv:2609.16921v1 (CC BY 4.0), distributed under the Creative Commons Attribution 4.0 International licence, as recorded in the arXiv licence metadata for this exact version and shown on the abstract page https://arxiv.org/abs/2609.16921v1. Full rights notice: licenses/arxiv-2609.16921-CC-BY-4.0.txt.\par\smallskip

\noindent\textit{Editorial scope.} Counted unit: the Theorem whose source label is \texttt{thm:01} in the article (source0.tex lines 731--780, source label \texttt{thm:01}), delivered together with the article's own complete original proof of it (source0.tex lines 782--1054, one \texttt{proof} environment of 273 lines). No mathematics is written, repaired or re-derived: statement and proof are the article's own text line for line. Required context retained uncounted: eq:10821 (lines 518--528); eq:0821 (lines 598--604); eq:082101 (lines 606--612). Every definition, display and label of those lines is retained unchanged. Omitted from this package and not redistributed: 1--517, 529--597, 605--605, 613--730, 781--781, 1055--1682. Nothing inside a retained excerpt is omitted or altered: every retained body is delivered exactly as the article's lines are written, so no omission receipt of any kind accompanies this package. Citations: the retained excerpts carry no citation command at all, so no bibliography entry of the article is transcribed and no reference is redistributed. Reference adaptation: none was needed and none was made; every reference inside the delivered unit resolves inside it, so no unresolved reference or citation is produced. Printed slips kept literal: no printed word, symbol, hypothesis or number of the counted statement or of its proof was altered. Layout and class: the article's own class is \texttt{IEEEtran} with packages including amssymb, amsthm, amsmath, latexsym, graphicx, mathrsfs, amsfonts, longtable, setspace, caption, rotating, ifsym, bbm, makecell; the wrapper is the lane's pinned \texttt{amsart} editorial wrapper at 10pt on A4, which restates the theorem-like environments the retained text needs and adds the article's own macro definitions that the retained text needs (\texttt{\textbackslash F}), with extra wrapper packages (bm, bbm, dsfont, xspace, cancel, mathtools). The delivered numbering is the unit's own. Assets: no retained span contains a figure environment, a graphics-inclusion command, a PDF-page inclusion, a table or a TikZ picture; no image file is copied into this package, the package declares no asset dependency and no image file is redistributed with it. Licence: version arXiv:2609.16921v1 of this article is distributed under the Creative Commons Attribution 4.0 International licence, as recorded in the arXiv licence metadata for this exact version and shown on the abstract page https://arxiv.org/abs/2609.16921v1. A full rights notice is delivered at licenses/arxiv-2609.16921-CC-BY-4.0.txt. Characters: every retained excerpt is pure ASCII, so the delivered engine, XeLaTeX with its default Latin Modern text font, reports no missing character.
\begin{lemma}\label{eq:10821}
Let $C$ and $D$ be linear codes of length $n$ and dimensions $k$ and $n-k$, respectively. Let $G_C$ and $G_D$ be generator matrices of $C$ and $D$. Then $(C,D)$ is an LCP if and only if
\[
G=
\begin{pmatrix}
G_C\\
G_D
\end{pmatrix}
\]
is nonsingular.
\end{lemma}

\begin{equation}\label{eq:0821}
\mathbf{C}_i=
\begin{cases}
\mathbf{V}_{i-1}+\eta_i\mathbf{V}_{k-1+t_i},&1\le i\le r,\\
\mathbf{V}_{i-1},&r<i\le k,
\end{cases}
\end{equation}

\begin{equation}\label{eq:082101}
\mathbf{D}_j=
\begin{cases}
\mathbf{V}_{j-1}+\delta_j\mathbf{V}_{n-k-1+\gamma_j},&1\le j\le s,\\
\mathbf{V}_{j-1},&s<j\le n-k.
\end{cases}
\end{equation}

\begin{theorem}\label{thm:01}
Let \(C\) and \(D\) be as above, and suppose
\[
\{n-k,\ldots,n-1\}
\subseteq
\{k-1+t_i:1\le i\le r\}
\cup
\{n-k-1+\gamma_j:1\le j\le s\}.
\]
Then \((C,D)\) is an LCP whenever one of the following conditions holds:

\begin{enumerate}
\renewcommand{\labelenumi}{(\roman{enumi})}

\item
\(0\le r<k\), \(s=k\), and
\(1\le t_i\le n-2k\) for \(1\le i\le r\).

\item
\(r=k\), all the exponents of the twist terms appearing in the
generator matrices of \(C\) and \(D\), namely
\[
k-1+t_1,\ldots,k-1+t_k, n-k-1+\gamma_1,\ldots,n-k-1+\gamma_s,
\]
are pairwise distinct, and
\(k-1+t_i\le n-k-1\) for \(1\le i\le s\), while
\(k-1+t_i\ge n-k\) for \(s+1\le i\le k\).

\item
\(0\le r<k\), \(s=k\),
\(1\le t_1<\cdots<t_r\), and
\(\gamma_i=i\) for \(1\le i\le k\).
Moreover, there exists \(j\) with \(1\le j\le r-1\) such that
\(
t_j\le n-2k<t_{j+1},
\)
and
$t_i=n-2k+i\text{ and }\eta_i\ne\delta_i$
for every \(j+1\le i\le r\).
\item
\(0\le r<k\), \(s=k\),
\(\gamma_i=i\) for \(1\le i\le k\), and $t_i=n-2k+i\text{ and }
\eta_i\ne\delta_i$ for every \(1\le i\le r\).

\item
\(r=k\), \(1\le t_1<\cdots<t_k\),
\(0\le s\le k\), \(n\le 2k-1+t_1\), and $\gamma_i=i\text{ and }
\eta_i\ne\delta_i$ for every \(1\le i\le s\).
\end{enumerate}
\end{theorem}

\begin{proof}
From Lemma~\ref{eq:10821}, it suffices to prove that
\[
G=
\begin{pmatrix}
G_C\\
G_D
\end{pmatrix}
\]
is nonsingular. Since
\[
V=
\begin{pmatrix}
\mathbf{V}_0\\
\vdots\\
\mathbf{V}_{n-1}
\end{pmatrix}
\]
is a nonsingular Vandermonde matrix, one may write \(G=AV\) for a uniquely determined coefficient matrix \(A\). Hence \(G\) is nonsingular if and only if \(A\) is nonsingular.

Equivalently, consider a linear relation
\begin{equation}\label{eq:ddw8}
\sum_{i=1}^{k}a_i\mathbf{C}_i
+
\sum_{j=1}^{n-k}b_j\mathbf{D}_j
=0,
\end{equation}
where \(\mathbf{C}_i\) and \(\mathbf{D}_j\) are given in
(\ref{eq:0821}) and (\ref{eq:082101}), respectively.
Because \(\mathbf{V}_0,\ldots,\mathbf{V}_{n-1}\) form a basis of \(\F_q^n\), the coefficient of each Vandermonde basis vector must vanish separately.

\medskip
\noindent\textbf{Case (i).}
Since \(s=k\) and \(\gamma_1,\ldots,\gamma_k\) are \(k\) pairwise distinct elements of \(\{1,\ldots,k\}\), we have
\(
\{\gamma_1,\ldots,\gamma_k\}=\{1,\ldots,k\}.
\)
Then the exponents of the twist terms in the generator matrix of \(D\) are exactly
\(
\{n-k,n-k+1,\ldots,n-1\},
\)
although not necessarily in this order.

On the other hand, by the assumption \(1\le t_i\le n-2k\),
\(
k-1+t_i\le n-k-1,
\)
so none of the exponents of the twist terms appearing in the generator matrix of \(C\) belongs to the set \(\{n-k,n-k+1,\ldots,n-1\}\). Comparing the coefficients corresponding to the exponents \(n-k,n-k+1,\ldots,n-1\) gives
$
b_j\delta_j=0,
$
where \(1\le j\le k\).
Since every \(\delta_j\) is nonzero, we have
\(
b_1=\cdots=b_k=0.
\)

We then consider the coefficients of the Vandermonde basis vectors
\(
\mathbf{V}_0,\ldots,\mathbf{V}_{k-1}.
\)
By the assumption \(1\le t_i\le n-2k\), all twist exponents appearing in \(G_C\) satisfy
\(
k\le k-1+t_i\le n-k-1,
\)
whereas all twist exponents appearing in \(G_D\) satisfy
\(
n-k\le n-k-1+\gamma_i\le n-1.
\)
Therefore, none of the twist terms in \(G_C\) or \(G_D\) contributes to the coefficients of
\(
\mathbf{V}_0,\ldots,\mathbf{V}_{k-1}.
\)
Hence, for each \(1\le i\le k\), the coefficient of \(\mathbf{V}_{i-1}\) in (\ref{eq:ddw8}) is
\(
a_i+b_i.
\)
Since \(b_1=\cdots=b_k=0\), we obtain
\(
a_1=\cdots=a_k=0.
\)

Finally, the remaining rows of \(G_D\) that do not contain twist terms correspond to the remaining Vandermonde basis vectors. Therefore,
\(
b_{k+1}=\cdots=b_{n-k}=0.
\)
Thus, (\ref{eq:ddw8}) has only the zero solution.

\medskip
\noindent\textbf{Case (ii).}
By the assumed exponent ranges, the first \(s\) twist terms in the
generator matrix of \(C\) have exponents in
\(
\{k,\ldots,n-k-1\},
\)
whereas the remaining \(k-s\) twist terms in \(C\) have exponents in
\(
\{n-k,\ldots,n-1\}.
\)

On the other hand, since \(1\le \gamma_j\le k\), every twist exponent
appearing in \(D\) satisfies
\(
n-k\le n-k-1+\gamma_j\le n-1,
\)
where \(1\le j\le s\).
Hence, among the twist terms appearing in \(C\) and \(D\), exactly
\(k-s\) twist terms from \(C\) and \(s\) twist terms from \(D\) have
exponents in
\(
\{n-k,\ldots,n-1\}.
\)
Since all twist exponents are pairwise distinct, these \(k\) exponents
occupy this \(k\)-element set exactly.

Moreover, each of these high-degree Vandermonde basis vectors occurs
in exactly one twist term among the generator rows of \(C\) and \(D\).
More precisely, for \(s+1\le i\le k\),
\(\mathbf{V}_{k-1+t_i}\) occurs only in the twist term of the
\(i\)-th row of \(C\), while for \(1\le j\le s\),
\(\mathbf{V}_{n-k-1+\gamma_j}\) occurs only in the twist term of the
\(j\)-th row of \(D\). Therefore, comparing the corresponding
coefficients in~(\ref{eq:ddw8}) gives
\[
\eta_i a_i=0\text{ and }\delta_j b_j=0,
\]
where \(s+1\le i\le k\),\text{ and }\(1\le j\le s\).
Since every \(\eta_i\) and \(\delta_j\) is nonzero, it follows that
\(a_{s+1}=\cdots=a_k=0\) and \(b_1=\cdots=b_s=0\).

For \(1\le i\le s\), neither the twist terms of \(C\) nor those of
\(D\) contribute to the coefficient of \(\mathbf{V}_{i-1}\).
Hence, comparison of the coefficient of \(\mathbf{V}_{i-1}\) gives
\(
a_i+b_i=0.
\)
Since \(b_i=0\), we obtain
\(
a_1=\cdots=a_s=0.
\)

Similarly, for \(s<i\le k\), comparison of the coefficient of
\(\mathbf{V}_{i-1}\) gives
\(
a_i+b_i=0.
\)
Since \(a_i=0\) for \(s<i\le k\), we obtain
\(
b_{s+1}=\cdots=b_k=0.
\)
Thus,
\(a_1=\cdots=a_k=0\) and \(b_1=\cdots=b_k=0\).
Substituting these equalities into~(\ref{eq:ddw8}), the linear relation
reduces to
\[
\sum_{j=k+1}^{n-k} b_j\mathbf{V}_{j-1}=0.
\]
Since
\(
\mathbf{V}_k,\mathbf{V}_{k+1},\ldots,
\mathbf{V}_{n-k-1}
\)
are linearly independent, it follows that
\(
b_{k+1}=\cdots=b_{n-k}=0.
\)
Consequently,~(\ref{eq:ddw8}) has only the zero solution.

\medskip
\noindent\textbf{Case (iii).}
For \(1\le i\le j\), the exponent of the twist term in \(C\) satisfies
\(
k-1+t_i\le n-k-1.
\)
Since \(\gamma_i=i\), the twist term in the \(i\)-th row of \(D\) has exponent
\(
n-k-1+i.
\)

For \(i\le j\), none of the twist terms in \(C\) has this exponent. Hence, comparison at \(\mathbf{V}_{n-k-1+i}\) gives
\(
b_i=0.
\)
Then the coefficient of \(\mathbf{V}_{i-1}\) gives
\(
a_i+b_i=0,
\)
and therefore \(a_i=0\).

For \(j+1\le i\le r\), the assumption \(t_i=n-2k+i\) gives
\(
k-1+t_i=n-k-1+i.
\)
Thus, the \(i\)-th twist terms in \(C\) and \(D\) have the same exponent. The corresponding coefficient equations are
\[
a_i+b_i=0 \text{ and } \eta_i a_i+\delta_i b_i=0.
\]
Equivalently,
\[
\begin{pmatrix}
1&1\\
\eta_i&\delta_i
\end{pmatrix}
\begin{pmatrix}
a_i\\
b_i
\end{pmatrix}
=0.
\]
Its determinant is \(\delta_i-\eta_i\ne0\), so \(a_i=b_i=0\).

For \(r<i\le k\), the row
\(
\mathbf{C}_i=\mathbf{V}_{i-1}
\)
contains no twist term, whereas the \(i\)-th twist term in \(D\) has exponent \(n-k-1+i\), which does not occur among the twist exponents of \(C\). Thus, comparison at \(\mathbf{V}_{n-k-1+i}\) gives \(b_i=0\), and comparison at \(\mathbf{V}_{i-1}\) gives \(a_i=0\).

Finally, the remaining untwisted rows of \(D\) correspond to the remaining Vandermonde basis vectors, so their coefficients also vanish. Therefore, (\ref{eq:ddw8}) has only the zero solution.

\medskip
\noindent\textbf{Case (iv).}
For every \(1\le i\le r\),
\(
k-1+t_i=n-k-1+i.
\)
Hence the twist terms in \(C\) and \(D\) contribute to the same Vandermonde basis vector \(\mathbf{V}_{n-k-1+i}\). The corresponding coefficient equations are identical to the \(2\times2\) system in Case (iii). Since \(\eta_i\ne\delta_i\), we obtain
\(
a_i=b_i=0.
\)
For \(r<i\le k\), the corresponding row of \(D\) contains a twist term with exponent \(n-k-1+i\), whereas no twist term in \(C\) has this exponent. Thus, comparison at \(\mathbf{V}_{n-k-1+i}\) gives \(b_i=0\), and comparison at \(\mathbf{V}_{i-1}\) gives \(a_i=0\).

Finally, all remaining coefficients vanish. Hence, (\ref{eq:ddw8}) has only the zero solution.

\medskip
\noindent\textbf{Case (v).}
Since
\(
1\le t_1<\cdots<t_k
\),\text{ and }
\(
n\le 2k-1+t_1,
\)
all \(k\) distinct exponents
\(
k-1+t_1,\ldots,k-1+t_k
\)
belong to the \(k\)-element set
\(
\{n-k,\ldots,n-1\}.
\)
Consequently,
\(
k-1+t_i=n-k+i-1,
\)
or equivalently,
\(
t_i=n-2k+i.
\)

For \(1\le i\le s\), the condition \(\gamma_i=i\) implies that the twist terms in \(C\) and \(D\) correspond to the same Vandermonde basis vector \(\mathbf{V}_{n-k+i-1}\). The corresponding coefficient equations form the same \(2\times2\) system as in Case (iii). Since
\(
\delta_i-\eta_i\ne0,
\) we obtain
\(
a_i=b_i=0.
\)

For \(s<i\le k\), the exponent \(n-k+i-1\) appears among the twist exponents of \(C\), but not among those of \(D\). Therefore, comparison at \(\mathbf{V}_{n-k+i-1}\) gives \(a_i=0\), and comparison at \(\mathbf{V}_{i-1}\) gives \(b_i=0\).

Finally, the remaining untwisted rows of \(D\) correspond to the remaining Vandermonde basis vectors, and their coefficients also vanish. Hence, (\ref{eq:ddw8}) has only the zero solution.

In every case, the stacked generator matrix is nonsingular, and therefore \((C,D)\) is an LCP.
\end{proof}

\end{document}
